\section{Relational Residual Bottlenecking}

\subsection{Setup}

Let $f_v$ and $f_t$ be frozen OpenCLIP image and text encoders. For an image $I$ or caption $T$, the frozen model produces normalized embeddings $\zbase^v=f_v(I)$ and $\zbase^t=f_t(T)$. The goal is to adjust the embedding used for contrastive comparison without discarding $\zbase$, since the pretrained path carries the broad retrieval and zero-shot behavior that we want to retain.

The method is designed for the dual-encoder setting. At inference time, images and captions must still be embeddable independently, and the final similarity remains a dot product between normalized vectors. No detector, parser, cross-attention reranker, scene graph, or candidate-pair module is added. This choice rules out some stronger reasoning machinery, but it keeps the evaluation focused on whether a pretrained embedding can be corrected rather than replaced.

RRB attaches a small trainable residual module to each branch:
\begin{equation}
    \zout^b = \mathrm{normalize}\left(\zbase^b + \alpha \drel^b\right), \qquad b \in \{v,t\}.
    \label{eq:rrb}
\end{equation}
Here $\drel^b$ is computed from CLIP token features through a parser-free multi-channel bottleneck and $\alpha$ is a fixed residual gate. The selected checkpoint uses $\alpha=0.05$. The residual therefore has to act as a correction to the original representation rather than a replacement representation.

The fixed gate is deliberately simple. A learned gate can increase flexibility, but the early R022 diagnostic run, omitted from the main tables for brevity, showed that more flexible residual control can still miss the COCO gate. A fixed gate is preferred because it maintains the architectural prior of a dominant pretrained path throughout training, preventing the residual branch from taking over the embedding space through an unconstrained mixing coefficient. The selected method fixes the residual scale and lets the objective decide when the residual direction is useful. This makes the resulting claims easier to audit: if guardrails fail, the failure cannot be attributed to an unbounded learned mixing coefficient silently replacing the frozen path.

\subsection{Multi-Channel Residual Interface}

For each branch, the module receives the frozen token sequence and induces $K$ latent channel states through learned queries and one self-attention block. The channel states are pooled and projected to the CLIP embedding dimension, producing $\drel^b$. We avoid the term ``semantic slots'' for the main method because later responsibility probes do not show semantic specialization. The bottleneck is used as an architectural constraint on the residual path, not as an interpretability claim.

The no-channel controls are important because they isolate this interface. R033 and R040 remove the multi-channel bottleneck and use a simpler residual adapter. Both remain relation-strong but damage retrieval and zero-shot behavior much more than RRB, suggesting that the bottlenecked residual interface is preservation-relevant even though its channels are not semantically interpretable.

The interface can be understood as a capacity and routing prior. The residual is not produced by one unconstrained multilayer projection from all token features to the final embedding. Instead, the token sequence must be summarized through a small number of channel states before it can affect the CLIP vector. This is not a guarantee of safety, and the seed results show that it is not enough by itself. It does, however, create a narrower path for the relation objective to act through. The contrast between R028 and the no-channel controls is consistent with this interpretation: similar or stronger relation accuracy can be achieved without the channel interface, but the resulting embedding drift is far more harmful to COCO and ImageNet.

This caution is consistent with recent analyses of ViT token aggregation, which show that global tokens can rely on background-dominated shortcut patches under coarse supervision~\cite{shi2026lastvit}. RRB does not attempt to fix dense-feature artifacts. The relevance is narrower: token summaries are not automatically benign, so the paper evaluates the final residual direction through retrieval and zero-shot guardrails instead of assuming that a small token-derived correction is safe.

The paper therefore separates three ideas that are easy to conflate. First, multiple channels are an architectural bottleneck. Second, the channels might become specialized in some future model. Third, the current probes do not show that specialization. Only the first statement is part of the method claim. The second is a hypothesis, and the third is a negative result reported in Section~4.

\subsection{Training Objective}

Training uses structured relation counterfactuals and anchor examples:
\begin{equation}
    \loss =
    \lambda_{\mathrm{rel}}\lintervene
    + \lambda_{\mathrm{anchor}}\lanchor
    + \lambda_{\Delta}\ldelta .
    \label{eq:loss}
\end{equation}
$\lintervene$ is a margin ranking loss over structured relation edits. Given a positive pair $(I,T^+)$ and a counterfactual caption $T^-$ that changes a relation, role, order, object, or attribute while preserving much of the surface vocabulary, the loss encourages $\simfn(\zout^v_I,\zout^t_{T^+})$ to exceed $\simfn(\zout^v_I,\zout^t_{T^-})$ by a margin. The same principle applies to branch-specific counterfactuals where available.

$\lanchor$ is feature mimicry on anchor pairs:
\begin{equation}
    \lanchor =
    \left\|\zout^v - \zbase^v\right\|_2^2 +
    \left\|\zout^t - \zbase^t\right\|_2^2 .
\end{equation}
It keeps the adapted embeddings close to the frozen OpenCLIP anchor on examples that should not require relational correction. $\ldelta$ penalizes residual magnitude,
\begin{equation}
    \ldelta = \left\|\drel^v\right\|_2^2 + \left\|\drel^t\right\|_2^2,
\end{equation}
which works with the fixed gate in Eq.~\ref{eq:rrb} to bound feature drift.

An earlier targeted relational loss attempted to localize edited spans to particular bottleneck channels. The final method removes that term. The ablations show that structured counterfactual supervision and residual anchoring are sufficient for the observed relation gain; the targeted loss is not necessary and should not be used as a mechanism claim.

\subsection{Structured Counterfactual Signal}

The relation loss depends on counterfactual captions whose surface form is close to the positive caption but whose relation, role, object, attribute, or order is changed. This structure is what distinguishes the method from generic hard-negative mining. Generic negatives can be visually or semantically difficult, but they need not isolate the binding error that CLIP tends to miss. In R024, generic hard negatives preserve broad OpenCLIP behavior almost perfectly, yet relation accuracy falls below the plain adapter baseline. That run is important because it rules out a simpler explanation: the selected model is not merely benefiting from more difficult negatives. It benefits from negatives that target the intended failure mode.

The training signal is also deliberately branch-compatible. Both image and text branches keep their frozen encoders, and both branches receive residual modules. The margin objective is evaluated through the final normalized embeddings, so the residual is rewarded only when it improves the contrastive comparison used by retrieval. This keeps the objective aligned with the deployed interface, although it also means that benchmark-specific counterfactual patterns can be learned without broader binding generalization. The Winoground and Flickr30k checks are included to expose that risk.

\subsection{Anchor Examples and Residual Regularization}

The anchor term uses examples that should not require a relational correction and asks the adapted embedding to mimic the frozen OpenCLIP embedding. This term is not a generic preference for small weights; it is a feature-space constraint at the point where downstream retrieval operates. The residual penalty is complementary. $\lanchor$ constrains the final normalized output on anchor examples, while $\ldelta$ constrains the branch residual itself. In the selected R028 configuration, $\lambda_{\mathrm{anchor}}=10$, $\lambda_{\Delta}=0.01$, and the anchor ratio is $0.75$.

The ablations show why both the anchor objective and the residual direction matter. R023 has low residual magnitude and high cosine to the frozen embedding, but it still misses the COCO gate. R028 keeps a similar small-residual regime and passes the pilot gate. R029 reintroduces the targeted loss and has a larger residual than R028 with no relation advantage. Anchor-20 runs reduce residual magnitude further but do not eliminate COCO seed fragility. The method should therefore be read as a bounded residual recipe, not as a theorem that smaller residuals are always safer.

\subsection{Model Selection Rule}

During experimentation, candidate checkpoints were evaluated with a pilot gate: relation accuracy had to exceed the plain adapter, COCO image-to-text R@1 had to be at least $0.70$, and ImageNet top-1 had to be at least $0.55$. This gate is not a universal deployment criterion. It is a project-level guard against selecting a model that wins the staged relation suite by destroying broad alignment. R028 is selected because it satisfies this gate while removing the unsupported targeted loss from the method. The seed studies then test how stable that selection is, rather than treating the selected checkpoint as sufficient evidence by itself.

For any practical use, this gate should be applied as a validation protocol rather than as a one-shot model choice. Multiple seeds should be trained under the same configuration; checkpoints that fail the relation, COCO, or ImageNet gates should be discarded; and the reported model should be accompanied by the seed pass rate. This protocol does not remove seed fragility, but it makes the fragility explicit and prevents relation-only selection from silently choosing an unsafe residual direction.
